\section{Eventual uniqueness among cubic maximizers}
\label{sec:four-coefficient-uniqueness}

\begin{corollary}
\label{cor:four-eventual-cubic-uniqueness}
There is an absolute integer $R$ such that, for every $r\ge R$, a connected cubic multigraph $G$ of rank $r$ satisfies
\[
 c(G)=\psi(r)
\]
if and only if $G$ is isomorphic to the Horak--Kim graph $MC_{r-1}$.
For odd $r$, this is the ring of $(r-1)/2$ square tiles.
For even $r$, it is the ring of $(r-2)/2$ square tiles and one single-rung tile.
The assertion concerns cubic graphs; it does not classify arbitrary subdivisions or other noncubic maximizers.
\end{corollary}

\begin{proof}
Increase $R$ until the eventual exact formula and every threshold in its proof apply.
Write $H_r=\psi(r)-2^{r-1}$ for the Horak--Kim excess.
Its two alternating increments are strictly positive.

The loop and high-girth deletion branches give excess at most $H_{r-1}<H_r$.
The bounded-girth branch has negative excess.
Thus every cubic maximizer has a cyclic cut of size at most four.
Cyclic cuts of size at most three give only $O(2^{r/2})$ excess, which is strictly smaller than $H_r$ for large $r$.

In the four-cut proof, every fixed-deficit branch lies below one half.
The aligned full-star branch again has only $O(2^{r/2})$ excess.
The only remaining case is an aligned pair of saturated cross poles with positive through counts.
It is a square and single-rung ring.

The fixed-rank optimization in Theorem~\ref{thm:four-coefficient-ring-maximum} is strict whenever the number of single-rung tiles exceeds its parity minimum.
The all-single-rung case is also strictly smaller.
Thus a maximizing ring has no single-rung tile at odd rank and exactly one at even rank.
Its cyclic tile order is unique up to rotation.

Finally, interface twists do not change its isomorphism type.
Swapping both strand labels in a tile transfers a twist from one adjacent interface to the other.
A square tile additionally permits swapping its left pair while fixing its right pair, since its internal graph is $K_{2,2}$.
Move all twists to an interface next to a square and remove the remaining twist by this last automorphism.
There is at least one square at the ranks under consideration.
The untwisted ring is exactly $MC_{r-1}$.
Conversely, that graph attains the eventual exact formula.
\end{proof}

The rank threshold is unspecified, as in the eventual exact-count theorem.
No uniqueness assertion is made at smaller ranks.
